%% ============================================================
%%  Chapter 3 — Methodology
%% ============================================================
\chapter{Methodology}
\label{ch:methodology}

This chapter describes the complete design of the HoneyIQ simulation
framework.
Section~\ref{sec:arch} presents the high-level architecture.
Section~\ref{sec:attacker} describes the attacker module, including the
session-coherent traffic-realism extension (\S\ref{sec:realism}).
Section~\ref{sec:threat} defines the composite threat-level metric and
reward function.
Section~\ref{sec:sedm} presents the Stage-Escalation Decision Matrix in
detail, including its fourth override rule.
Section~\ref{sec:dqn_baseline} describes the DQN baseline.
Section~\ref{sec:env} documents the Gymnasium environment.
Section~\ref{sec:classifier} covers the Random Forest classifier.
Section~\ref{sec:eval_infra} describes the evaluation infrastructure.
Section~\ref{sec:escalation_tracking} introduces the severity-weighted
escalation-tracking alternative.
Section~\ref{sec:dynamic_response} presents the two auditable
dynamic-response mechanisms.

%% ── 3.1 System Architecture ──────────────────────────────────
\section{System Architecture}
\label{sec:arch}

HoneyIQ is a closed-loop simulation in which an \emph{Attacker} and a
\emph{Defender} interact through a Gymnasium-compatible environment
(\texttt{CyberSecurityEnv}).
Figure~\ref{fig:architecture_text} shows the component layout and data flow.

\begin{figure}[ht]
  \centering
  \begin{tabular}{|p{0.85\textwidth}|}
    \hline
    \vspace{0.3cm}
    \textbf{CyberSecurityEnv} (gymnasium.Env wrapper)
    \begin{description}
      \item[Attacker:] TransitionModel (Markov chains) $\rightarrow$
            session-coherent feature simulation (\S\ref{sec:realism})
            $\rightarrow$ $(a_t, k_t, \mathbf{x}_t)$
      \item[Threat:] compute\_threat\_level$(a_t, k_t, r_t, n_t) \rightarrow T \in [0,1]$
      \item[State:] $s_t \in \mathbb{R}^{24}$ (one-hot + scalars)
      \item[Defender:] AttackClassifier (RF) $+$ MatrixPolicy (SEDM) $\rightarrow$ action $\in \{0,\ldots,4\}$
      \item[Reward:] compute\_reward$(a_t, T, k_t, \text{attack\_type}) \rightarrow r_t$
    \end{description}
    \vspace{0.1cm}
    \textbf{MetricsCollector}: StepRecord $\rightarrow$ EpisodeRecord $\rightarrow$ CSV + plots
    \vspace{0.3cm} \\
    \hline
  \end{tabular}
  \caption{High-level architecture of the HoneyIQ simulation framework.}
  \label{fig:architecture_text}
\end{figure}

At each discrete time step $t$, the following sequence of operations is
executed:
%
\begin{enumerate}
  \item The attacker advances its Markov chains to produce a new
        $(a_t, k_t)$ pair and samples a 15-dimensional feature vector
        $\mathbf{x}_t$, consistent with that session's persistent
        intensity/persona profile (\S\ref{sec:realism}).
  \item The environment updates the escalation-rate signal(s)
        (\S\ref{sec:escalation_tracking}), computes the threat level $T_t$,
        and constructs the 24-dimensional state vector $s_t$.
  \item The defender observes $s_t$ (and, in the live-pipeline deployment
        described in \S\ref{sec:dynamic_response}, a cross-session
        reputation score) and selects a honeypot action.
  \item The reward $r_t$ is computed and returned to the defender.
  \item Metrics are recorded.
\end{enumerate}

This one-step interaction repeats for a configurable number of steps
(200 in the original evaluation, 500 in the revised primary evaluation and
in DQN training — see Chapter~\ref{ch:results}) before the episode terminates.

%% ── 3.2 Attacker Module ──────────────────────────────────────
\section{Attacker Module}
\label{sec:attacker}

\subsection{Attack Type and Kill Chain Enumerations}

The discrete state spaces of the attacker are defined by two \texttt{IntEnum}
classes:

\paragraph{AttackType.}
Ten categories inspired by the UNSW-NB15 classification taxonomy:
NORMAL (0), RECONNAISSANCE (1), ANALYSIS (2), FUZZERS (3), GENERIC (4),
EXPLOITS (5), SHELLCODE (6), BACKDOORS (7), DOS (8), WORMS (9).
Each category is assigned a severity weight $s_i \in [0, 1]$
(Table~\ref{tab:attack_severity_meth}).

\begin{table}[ht]
  \centering
  \caption{Attack types, severity weights, primary kill-chain stages,
           and kill-chain stage weights.}
  \label{tab:attack_severity_meth}
  \begin{tabular}{clccc}
    \toprule
    Index & Attack Type & Severity $s_i$ & Primary Stage & Stage Weight $w_k$ \\
    \midrule
    0 & NORMAL        & 0.00 & Reconnaissance (0) & 0.10 \\
    1 & RECONNAISSANCE & 0.20 & Reconnaissance (0) & 0.10 \\
    2 & ANALYSIS      & 0.25 & Weaponization (1)  & 0.20 \\
    3 & FUZZERS       & 0.35 & Delivery (2)       & 0.35 \\
    4 & GENERIC       & 0.40 & Delivery (2)       & 0.35 \\
    5 & EXPLOITS      & 0.70 & Exploitation (3)   & 0.55 \\
    6 & SHELLCODE     & 0.75 & Exploitation (3)   & 0.55 \\
    7 & BACKDOORS     & 0.80 & Installation (4)   & 0.70 \\
    8 & DOS           & 0.85 & Actions on Obj (6) & 1.00 \\
    9 & WORMS         & 0.90 & Cmd \& Ctrl (5)    & 0.85 \\
    \bottomrule
  \end{tabular}
\end{table}

\paragraph{KillChainStage.}
Seven stages mapped to the Lockheed Martin Kill Chain:
RECONNAISSANCE (0), WEAPONIZATION (1), DELIVERY (2), EXPLOITATION (3),
INSTALLATION (4), COMMAND\_AND\_CTRL (5), ACTIONS\_ON\_OBJ (6).
Stage weights $w_k$ increase monotonically from 0.10 to 1.00, reflecting
the escalating danger of later kill-chain activity.

\subsection{Markov Transition Model}
\label{sec:transition}

Two row-stochastic base transition matrices are maintained:
$\mathbf{T}_A \in \mathbb{R}^{10 \times 10}$ (attack types) and
$\mathbf{T}_K \in \mathbb{R}^{7 \times 7}$ (kill chain stages).
At each step the attacker samples:
%
\begin{align}
  a_{t+1} &\sim \text{Categorical}\bigl(\mathbf{T}_A^{(\pi)}[a_t, \cdot]\bigr), \\
  \tilde{k}_{t+1} &\sim \text{Categorical}\bigl(\mathbf{T}_K^{(\pi)}[k_t, \cdot]\bigr).
\end{align}
%
A kill-chain floor constraint prevents unrealistic backward jumps:
%
\begin{equation}
  k_{t+1} = \max\!\bigl(\tilde{k}_{t+1},\;
                          \text{primary\_stage}(a_{t+1}) - 1\bigr).
  \label{eq:floor}
\end{equation}
%
This ensures that attack types associated with late kill chain stages cannot
appear at very early stages, maintaining temporal consistency.

\subsubsection{Intent Profiles}
\label{sec:intents}

Four intent profiles $\pi \in \{\text{STEALTHY}, \text{AGGRESSIVE},
\text{TARGETED}, \text{OPPORTUNISTIC}\}$ are encoded as element-wise
modifier matrices applied before row-renormalisation
(Table~\ref{tab:intents_meth}).

\begin{table}[ht]
  \centering
  \caption{Attacker intent profiles, behavioural characteristics, and
           primary attack preferences.}
  \label{tab:intents_meth}
  \begin{tabular}{lll}
    \toprule
    Intent & Attack Preference & Kill Chain Speed \\
    \midrule
    STEALTHY      & Recon, Analysis, Backdoors    & Slow; high self-loop probability \\
    AGGRESSIVE    & DoS, Worms, Exploits          & Fast; strong forward bias \\
    TARGETED      & Exploits $\to$ Shellcode $\to$ Backdoors & Direct exploitation chain \\
    OPPORTUNISTIC & Generic, Fuzzers (scattered)  & Moderate forward bias \\
    \bottomrule
  \end{tabular}
\end{table}

\textbf{STEALTHY}: the modifier matrix amplifies the self-loop probability
for RECONNAISSANCE and ANALYSIS, and increases the transition probability
from ANALYSIS to BACKDOORS.
Stage transitions are biased towards self-loops (the attacker dwells in each
stage longer), producing long-duration, low-severity campaigns.

\textbf{AGGRESSIVE}: forward transitions are amplified for DOS, WORMS, and
EXPLOITS.
Stage transitions are biased strongly forward, producing rapid escalation
through the kill chain.
The result is short, high-intensity campaigns that quickly reach ACTIONS\_ON\_OBJ.

\textbf{TARGETED}: transitions from EXPLOITS to SHELLCODE, SHELLCODE to
BACKDOORS, and BACKDOORS to WORMS are amplified, encoding a specific attack
chain.
The kill chain transitions are biased towards a fast path from EXPLOITATION
through INSTALLATION to COMMAND\_AND\_CTRL.

\textbf{OPPORTUNISTIC}: transitions to FUZZERS and GENERIC are amplified,
producing scattered attack-type sequences.
Stage transitions maintain moderate forward bias, resulting in a campaign that
explores multiple attack types without strong directionality.

\subsection{Network Feature Simulation}
\label{sec:features}

At each step the attacker samples a 15-dimensional feature vector
$\mathbf{x}_t \in \mathbb{R}^{15}$ from parametric distributions conditioned
on the current attack type.
The fifteen features are:
\texttt{dur}, \texttt{sbytes}, \texttt{dbytes}, \texttt{sttl},
\texttt{dttl}, \texttt{sloss}, \texttt{dloss}, \texttt{sload},
\texttt{dload}, \texttt{spkts}, \texttt{dpkts}, \texttt{swin},
\texttt{dwin}, \texttt{ct\_srv\_src}, \texttt{ct\_dst\_ltm}.

Selected distributions illustrate the per-attack signature design:

\begin{itemize}
  \item \textbf{RECONNAISSANCE}: \texttt{dur}~$\sim$~Exponential(rate=5),
        \texttt{ct\_dst\_ltm}~$\sim$~Poisson(30), reflecting rapid
        port-scanning activity.
  \item \textbf{DoS}: \texttt{sload}~$\sim$~Uniform(50\,000, 500\,000)\,bps,
        \texttt{dur}~$\sim$~Exponential(rate=10), capturing the high-bandwidth,
        short-burst pattern of flooding attacks.
  \item \textbf{BACKDOORS}: \texttt{dur}~$\sim$~Uniform(10, 3600)\,s,
        \texttt{sload}~$\sim$~Uniform(50, 500)\,bps, modelling long-lived,
        low-bandwidth covert channels.
  \item \textbf{WORMS}: \texttt{ct\_dst\_ltm}~$\sim$~Poisson(40),
        \texttt{spkts}~$\sim$~Poisson(200), reflecting spreading behaviour.
  \item \textbf{EXPLOITS}: \texttt{sbytes}~$\sim$~Uniform(1000, 50\,000),
        \texttt{sload}~$\sim$~Uniform(10\,000, 100\,000)\,bps, capturing
        the payload-heavy pattern of exploit delivery.
\end{itemize}

The 15 features are passed to the Random Forest classifier for attack-type
identification.
They are not directly included in the DQN state vector, which uses
ground-truth labels.

\textit{Data format generality note:} although the 15-field schema above is
named after UNSW-NB15 for familiarity, neither the classifier nor the SEDM
has a code-level dependency on UNSW-NB15 or NSL-KDD specifically — both
operate on a plain \texttt{dict[str, float]} matching this schema, which maps
onto other flow-export formats (NetFlow v9/IPFIX, Zeek \texttt{conn.log},
CICFlowMeter). See Appendix~\ref{app:hyperparams} for the parameter
listing.

\subsection{Session-Coherent Traffic Realism}
\label{sec:realism}

The feature simulator described above draws every one of the fifteen
features independently, at every step, from the attack type's static
distribution. Two consequences follow directly from that design, and
motivate the extension presented here. First, a simulated session has no
persistent ``character''\,: nothing distinguishes a session generated by a
consistently high-volume source from one generated by a consistently
low-volume source, because every step is an independent draw. Second, all
\texttt{NORMAL}-labelled traffic — the entirety of the benign class — is
drawn from a single fixed distribution, so the simulator cannot represent
the fact that real benign traffic on a network arrives from qualitatively
different kinds of source (an interactive user, an automated crawler, a
periodic health check).

\paragraph{Per-session intensity.}
\texttt{Attacker} now draws a persistent intensity scalar once per episode,
%
\begin{equation}
  \iota \sim \text{Lognormal}(0,\, \sigma_\iota), \qquad \sigma_\iota = 0.35,
  \label{eq:intensity}
\end{equation}
%
applied multiplicatively to the eight volume-shaped features
(\texttt{sbytes}, \texttt{dbytes}, \texttt{sload}, \texttt{dload},
\texttt{spkts}, \texttt{dpkts}, \texttt{ct\_srv\_src}, \texttt{ct\_dst\_ltm})
for the remainder of that episode. The remaining seven features
(\texttt{dur}, \texttt{sttl}, \texttt{dttl}, \texttt{sloss}, \texttt{dloss},
\texttt{swin}, \texttt{dwin}) are left unscaled: scaling a TCP window size or
a time-to-live value by a session-level ``intensity'' factor has no physical
interpretation, whereas byte counts, packet counts, and connection counts
plausibly do scale together for a single, consistently-behaving source. The
intensity draw is repeated on every \texttt{reset()} (after the underlying
random number generator is re-seeded, if a seed was supplied), so under a
fixed seed the sequence of per-episode intensities is itself reproducible.

\paragraph{Benign traffic personas.}
For \texttt{NORMAL}-labelled traffic specifically, the simulator additionally
draws a categorical persona once per episode:
%
\begin{equation}
  \text{persona} \sim \text{Categorical}\bigl(\{\texttt{casual\_user}: 0.70,\;
  \texttt{crawler}: 0.20,\; \texttt{monitoring\_probe}: 0.10\}\bigr).
\end{equation}
%
\texttt{casual\_user} reuses the original NORMAL distribution unchanged.
\texttt{crawler} represents a short-duration, high-request-rate,
single-service source (high \texttt{ct\_srv\_src}, low \texttt{ct\_dst\_ltm},
near-zero loss). \texttt{monitoring\_probe} represents a very short, very
small, many-destination source (high \texttt{ct\_dst\_ltm}, minimal
payload). Persona selection affects only which of the fifteen feature
distributions is sampled; it has no effect on \texttt{attack\_type},
\texttt{kill\_chain\_stage}, or any downstream one-hot encoding —
\texttt{NORMAL} remains \texttt{NORMAL} regardless of persona.

\paragraph{Independent-sample override.}
Two call sites in the system require many \emph{independent} samples of a
single attack type rather than one session's worth of session-coherent
samples: \texttt{AttackClassifier}'s training-data generator
(\S\ref{sec:classifier}) and the feature-distribution diagnostic
visualisation. Both were updated to draw a fresh intensity (and, for
\texttt{NORMAL}, a fresh persona) per sample from a local random number
generator, passed explicitly to the feature simulator via keyword-only
override parameters that default to the calling session's own profile.
Without this explicit override, every ``independent'' training sample of a
given class would silently share one session's intensity and persona,
understating the feature variance the classifier is trained to discriminate
against — a bias that acts specifically to make classification appear
\emph{easier} than it should, since it manufactures artificial within-class
consistency that would not exist in a population of genuinely distinct
sources.

\paragraph{Randomised event payloads.}
Independently of the feature simulator, the OpenCanary-integration event
generator's payload templates — previously a small, fixed set of literal
example credentials, paths, and counts per scenario — are now resolved from
small wordlists and jitter ranges at generation time, so that repeated
invocations of the same named scenario (e.g.\ \texttt{ssh\_brute}) produce
varied, plausible-but-not-identical payloads rather than one of a handful of
exact literal strings. This affects only the descriptive \texttt{logdata}
payload of a generated event; the kill-chain classification pipeline
(\S\ref{sec:classifier}, \S\ref{sec:dynamic_response}) reads only the
numeric logtype code and is unaffected.

The empirical effect of these extensions on classifier accuracy and
downstream SEDM metrics is reported in Chapter~\ref{ch:results}.

%% ── 3.3 Composite Threat Level and Reward ────────────────────
\section{Composite Threat Level and Reward Function}
\label{sec:threat}

\subsection{Composite Threat Level}
\label{sec:threat_level_def}

The threat level $T_t \in [0, 1]$ aggregates four signals into a single
normalised index:
%
\begin{equation}
  T_t = 0.45 \cdot s_{a_t}
      + 0.35 \cdot w_{k_t}
      + 0.15 \cdot r_t
      + 0.05 \cdot \min\!\left(1,\, \frac{n_t}{100}\right),
  \label{eq:threat}
\end{equation}
%
where:
\begin{itemize}
  \item $s_{a_t} \in [0, 0.90]$ is the attack-type severity weight
        (Table~\ref{tab:attack_severity_meth}).
  \item $w_{k_t} \in [0.10, 1.00]$ is the kill-chain stage weight.
  \item $r_t \in [0, 1]$ is the \emph{escalation rate}, defined and — as of
        this version — computed in one of two ways
        (\S\ref{sec:escalation_tracking}). Unless stated otherwise, all
        results in this thesis prior to the extension chapters use the
        original window-based definition: the fraction of the last 20
        steps that contained attack traffic.
  \item $n_t$ is the cumulative attack count; the capped normalisation
        $\min(1, n_t / 100)$ provides a measure of cumulative campaign
        pressure that saturates after 100 attacks.
\end{itemize}

The weight assignment reflects a domain-informed prioritisation.
Attack severity (45\%) is the dominant signal because the attack type
directly indicates the attacker's capability and intent.
Kill chain stage (35\%) captures the attacker's progress through the
campaign, which determines the urgency of response.
Escalation rate (15\%) and cumulative count (5\%) provide temporal context
that distinguishes an isolated event from an ongoing campaign.

\paragraph{Threat bands.}
The threat level is partitioned into five bands used by the reward function:

\begin{center}
\begin{tabular}{lc}
  \toprule
  Band & Range \\
  \midrule
  Benign & $T < 0.15$ \\
  Low    & $0.15 \leq T < 0.35$ \\
  Medium & $0.35 \leq T < 0.55$ \\
  High   & $0.55 \leq T < 0.75$ \\
  Critical & $T \geq 0.75$ \\
  \bottomrule
\end{tabular}
\end{center}

\subsection{Reward Function}
\label{sec:reward}

The reward function $R(a, T, k, \text{attack\_type})$ encodes domain
knowledge as a structured signal with three layers.

\paragraph{Layer 1: Base matrix.}
The base reward is looked up from a $5 \times 5$ action-by-threat-band matrix
(Table~\ref{tab:reward_matrix_meth}).

\begin{table}[ht]
  \centering
  \caption{Base reward matrix (action $\times$ threat band).
           Values encode the desirability of each action--threat pairing.}
  \label{tab:reward_matrix_meth}
  \begin{tabular}{lccccc}
    \toprule
    Action & Benign & Low & Medium & High & Critical \\
    \midrule
    ALLOW & $+1.0$ & $+0.5$ & $-1.0$ & $-3.0$ & $-6.0$ \\
    LOG   & $+0.2$ & $+1.5$ & $+2.0$ & $+1.0$ & $-1.0$ \\
    TROLL & $-1.0$ & $+1.0$ & $+3.0$ & $+2.5$ & $+0.5$ \\
    BLOCK & $-2.0$ & $-0.5$ & $+1.5$ & $+3.5$ & $+5.0$ \\
    ALERT & $-3.0$ & $-1.0$ & $+0.5$ & $+2.0$ & $+6.0$ \\
    \bottomrule
  \end{tabular}
\end{table}

The matrix encodes graduated proportional responses: ALLOW is rewarded for
benign traffic but incurs increasing penalties as threat level rises.
ALERT receives the highest reward for critical threats ($+6.0$) and the
largest penalty for benign traffic ($-3.0$), reflecting its asymmetric cost
structure.
TROLL peaks at Medium threat (+3.0), incentivising engagement with
medium-severity attacks for intelligence gathering.

\paragraph{Layer 2: Late kill-chain amplifier.}
At kill chain stages Installation (4), Command \& Control (5), and
Actions on Objectives (6), negative rewards are amplified by a factor of
1.5$\times$:
%
\begin{equation}
  r \leftarrow 1.5 \cdot r \quad \text{if } r < 0 \text{ and } k \geq 4.
\end{equation}
%
This modifier encodes the escalating cost of defensive errors at advanced
campaign stages, where allowing or under-responding to an attack is
substantially more damaging than at early stages.

\paragraph{Layer 3: Attack-type bonuses.}
Specific action--attack-type combinations receive domain-knowledge bonuses:
%
\begin{itemize}
  \item TROLL on BACKDOORS, SHELLCODE, or WORMS: $+0.8$.
        These persistent attack types offer the highest intelligence value
        when the defender engages them for extended periods.
  \item BLOCK on WORMS: $+1.0$.
        Rapid containment of spreading malware is particularly valuable.
  \item LOG on RECONNAISSANCE: $+0.5$.
        Silent logging of reconnaissance provides intelligence without
        alerting the attacker.
  \item ALLOW on confirmed non-attack: $+0.5$.
        Correct passthrough of benign traffic is explicitly rewarded to
        discourage over-aggressive responses.
\end{itemize}

The reward function is unchanged in this version. It continues to serve
episode-level bookkeeping and DQN training (\S\ref{sec:dqn_baseline}); the
SEDM, being a fixed lookup policy, never optimises against it, so none of
the extensions in \S\ref{sec:escalation_tracking}--\S\ref{sec:dynamic_response}
alter the reward function itself.

%% ── 3.4 Stage-Escalation Decision Matrix ─────────────────────
\section{Stage-Escalation Decision Matrix (SEDM)}
\label{sec:sedm}

The SEDM is a transparent, deterministic policy that maps the observable
environment state to a honeypot action through a five-step procedure,
extended in this version with a fourth override rule.

\subsection{Design Rationale}

The SEDM is motivated by three observations:

\begin{enumerate}
  \item The Markov chain structure provides \emph{exact} escalation
        probabilities that encode the attacker's likely next move without
        requiring trajectory data.
  \item Kill chain stage is the single most informative variable for
        honeypot response: the appropriate action is qualitatively
        different at Reconnaissance (observe passively) versus
        Actions on Objectives (escalate immediately).
  \item Operational security practitioners require policies whose reasoning
        can be audited, documented, and modified by hand.
        A $7 \times 3$ matrix with a small number of override rules
        satisfies this requirement; a 256-node neural network does not.
\end{enumerate}

\subsection{Algorithm}
\label{sec:sedm_alg}

\paragraph{Step 1 — Escalation risk.}
Query the intent-specific Markov transition model for the probability of
advancing to a strictly higher kill chain stage from the current stage $k$:
%
\begin{equation}
  \rho(k, \pi) = \sum_{k' > k} T^{(\pi)}_{kk'}.
  \label{eq:esc_risk_meth}
\end{equation}
%
This quantity is computed analytically from the stored transition matrix.
Under the STEALTHY intent, stage transitions are biased towards self-loops,
producing low escalation risk values.
Under the AGGRESSIVE intent, strong forward transitions produce high
escalation risk.

\paragraph{Step 2 — Band classification.}
Discretise $\rho$ into three bands:
%
\begin{equation}
  \text{band} =
  \begin{cases}
    0 \; (\text{Low})    & \rho < 0.35 \\
    1 \; (\text{Medium}) & 0.35 \leq \rho < 0.65 \\
    2 \; (\text{High})   & \rho \geq 0.65.
  \end{cases}
\end{equation}
%
The thresholds 0.35 and 0.65 partition the escalation risk into three
roughly equal tertiles.
The Low band covers stages with strong self-loops (primarily Reconnaissance
under Stealthy intent); High covers stages with strong forward momentum
(primarily Delivery through Installation under Aggressive intent). These two
thresholds are deliberately excluded from the dynamic-adjustment mechanisms
introduced in \S\ref{sec:dynamic_response}, for a reason made precise there:
$\rho(k, \pi)$ has no live observational component, so there is no honest
runtime signal against which to adapt them.

\paragraph{Step 3 — Matrix lookup.}
Read the base action from the $7 \times 3$ SEDM
(Table~\ref{tab:sedm_matrix}):
%
\begin{equation}
  a_{\text{base}} = \text{SEDM}[k][\text{band}].
\end{equation}

\begin{table}[ht]
  \centering
  \caption{Stage-Escalation Decision Matrix.
           Rows are kill chain stages; columns are escalation risk bands.}
  \label{tab:sedm_matrix}
  \begin{tabular}{lccc}
    \toprule
    Stage & Low ($\rho < 0.35$) & Medium ($0.35 \leq \rho < 0.65$) & High ($\rho \geq 0.65$) \\
    \midrule
    RECONNAISSANCE   & ALLOW & LOG   & LOG   \\
    WEAPONIZATION    & LOG   & LOG   & TROLL \\
    DELIVERY         & LOG   & TROLL & TROLL \\
    EXPLOITATION     & TROLL & BLOCK & BLOCK \\
    INSTALLATION     & BLOCK & BLOCK & ALERT \\
    COMMAND\_\&\_CTRL & BLOCK & ALERT & ALERT \\
    ACTIONS\_ON\_OBJ & ALERT & ALERT & ALERT \\
    \bottomrule
  \end{tabular}
\end{table}

The matrix encodes three design principles:
%
\begin{itemize}
  \item \textbf{Proportionality}: responses escalate with both kill chain
        stage and escalation risk. No ALERT is issued before Exploitation
        even under high escalation risk; no ALLOW is issued after
        Weaponization.
  \item \textbf{Intelligence over containment at low stages}: LOG and TROLL
        are preferred over BLOCK at early stages to maximise intelligence
        before committing to containment.
  \item \textbf{Decisive containment at late stages}: ALERT is the
        dominant action from Installation onwards under medium or high
        escalation risk, reflecting the low tolerance for delay at advanced
        kill chain positions.
\end{itemize}

\paragraph{Step 4 — Override rules.}
Four override rules are applied, in the following priority order, after the
matrix lookup:
%
\begin{description}
  \item[R4 (Cross-session reputation):]
        If the requesting source's cross-session reputation score
        (\S\ref{sec:reputation}) is at or above \texttt{REPUTATION\_THRESHOLD}
        $= 0.60$, upgrade the action by one level and stop — this rule is
        checked \emph{before} R1. Because reputation is the only override
        input that concerns a source's \emph{history} rather than its
        current traffic, evaluating it first ensures a previously-established
        repeat offender cannot regain unconditional trust by sending a single
        innocuous-looking event; \S\ref{sec:reputation} discusses the
        trade-off this ordering implies.
  \item[R1 (Normal traffic):]
        If \texttt{AttackType} = NORMAL and R4 did not trigger, return ALLOW
        regardless of the matrix result. This ensures that confirmed benign
        traffic from a source with no accumulated reputation is never
        blocked.
  \item[R2 (High-impact attack types):]
        If \texttt{AttackType} $\in$ \{DOS, WORMS\} and neither R4 nor R1
        triggered, upgrade the action by one level in the sequence
        ALLOW $\to$ LOG $\to$ TROLL $\to$ BLOCK $\to$ ALERT.
        These attack types cause rapid, widespread damage and warrant
        an immediately more aggressive response.
  \item[R3 (High attack frequency):]
        If the escalation rate $r_t$ exceeds the (possibly
        adaptively-tuned, \S\ref{sec:threshold_controller}) rate threshold
        and no earlier rule triggered, upgrade the action by one level.
        A sliding window (or, under EMA tracking, a smoothly-decaying
        signal — \S\ref{sec:escalation_tracking}) showing sustained attack
        activity indicates an ongoing campaign that requires escalated
        containment.
\end{description}

Rules are evaluated in strict priority order; the first rule whose condition
holds determines the final action, and no step applies more than one
upgrade. This precedence extends, without otherwise altering, the
three-rule sequence (R1, R2, R3) used prior to this version; a system with
\texttt{reputation} always equal to zero — the value used throughout
Chapters~\ref{ch:results}--\ref{ch:discussion} except where stated otherwise —
reproduces the original three-rule behaviour exactly, since R4's condition
can never hold.

\paragraph{Step 5 — Composite risk score.}
A composite risk score $\rho_c \in [0, 1]$ is computed for logging and
analysis.
It does not affect the action selection:
%
\begin{equation}
  \rho_c = 0.35 \cdot w_k
         + 0.35 \cdot \rho(k, \pi)
         + 0.15 \cdot s_a
         + 0.15 \cdot r_t.
  \label{eq:composite_risk}
\end{equation}

\subsection{Intent-Awareness}

The SEDM is intent-aware through the escalation risk computation
(Equation~\ref{eq:esc_risk_meth}).
The base matrix $\text{SEDM}[k][\text{band}]$ is identical across all
intents, but the band assigned to a given stage differs:
%
\begin{itemize}
  \item Under STEALTHY, stage transitions are slow.
        RECONNAISSANCE has low escalation risk ($\rho < 0.35$) $\to$ ALLOW.
        The defender correctly passes low-threat reconnaissance traffic.
  \item Under AGGRESSIVE, stage transitions are fast.
        Even RECONNAISSANCE may have medium escalation risk ($\rho \geq 0.35$)
        $\to$ LOG, reflecting the attacker's tendency to escalate quickly.
  \item Under TARGETED, the direct exploitation path produces consistently
        high escalation risk from EXPLOITATION onwards $\to$ BLOCK/ALERT
        dominates.
\end{itemize}

The escalation risk is the key mechanism through which the SEDM adapts to
different attacker behaviours without any learning or parameter updates.

\subsection{Worked Example}

Consider an observation with:
stage = EXPLOITATION (3), attack type = EXPLOITS, escalation rate = 0.65,
intent = AGGRESSIVE, reputation = 0.0.

\begin{enumerate}
  \item \textbf{Escalation risk}: Under AGGRESSIVE intent,
        $T^{(\text{AGG})}_{3,4} + T^{(\text{AGG})}_{3,5} + T^{(\text{AGG})}_{3,6} = 0.72$
        (high forward bias in the transition matrix).
  \item \textbf{Band}: $\rho = 0.72 \geq 0.65 \Rightarrow$ band = High.
  \item \textbf{Matrix lookup}: SEDM[3][2] = BLOCK.
  \item \textbf{Override check}: reputation $0.0 < 0.60$ (R4 not triggered);
        EXPLOITS $\notin$ \{DOS, WORMS\} (R2 not triggered);
        escalation rate 0.65 $\leq$ 0.80 (R3 not triggered).
  \item \textbf{Final action}: BLOCK.
  \item \textbf{Composite risk}:
        $\rho_c = 0.35 \times 0.55 + 0.35 \times 0.72 + 0.15 \times 0.70 + 0.15 \times 0.65 = 0.654$.
\end{enumerate}

Now suppose the same observation instead arrives from a source with
reputation $= 0.75$ (e.g.\ the source has accumulated four or more prior
EXPLOITS-severity offenses — see Chapter~\ref{ch:results}). R4's condition
$0.75 \geq 0.60$ now holds and is checked first: the final action becomes
the one-level upgrade of the matrix's base lookup, ALERT, regardless of what
R1--R3 would otherwise have produced.

%% ── 3.5 DQN Baseline ─────────────────────────────────────────
\section{DQN Baseline}
\label{sec:dqn_baseline}

The DQN baseline provides a learned policy for comparison with the SEDM.
Its architecture and training hyperparameters are summarised in
Table~\ref{tab:hyperparams}.

\begin{table}[ht]
  \centering
  \caption{DQN hyperparameters used in all training experiments.}
  \label{tab:hyperparams}
  \begin{tabular}{ll}
    \toprule
    Parameter & Value \\
    \midrule
    State dimension & 24 \\
    Action space & 5 (\texttt{HoneypotAction}) \\
    Hidden layers & [256, 128, 64] \\
    Normalisation per layer & \texttt{LayerNorm} \\
    Activation & \texttt{ReLU} \\
    Weight initialisation & Kaiming uniform \\
    Replay buffer capacity & 15\,000 transitions \\
    Mini-batch size & 64 \\
    Discount factor $\gamma$ & 0.99 \\
    Learning rate & $10^{-3}$ (Adam) \\
    Target network update interval & 150 steps (hard copy) \\
    $\varepsilon$ initial / final / decay & 1.0 / 0.05 / 0.997 \\
    Gradient clip $\ell_2$ norm & 10.0 \\
    Loss function & Huber (SmoothL1) \\
    Training episodes & 300 \\
    Steps per episode & 500 \\
    Training intent & OPPORTUNISTIC \\
    \bottomrule
  \end{tabular}
\end{table}

\subsection{Training Protocol}

Training proceeds for 300 episodes of 500 steps each, using the
OPPORTUNISTIC attacker intent.
OPPORTUNISTIC was chosen as the training distribution because its scattered
attack-type sequence provides broad coverage of the state space, encouraging
the policy network to learn a generalised Q-function rather than one
optimised for a specific escalation pattern.

At each step:
%
\begin{enumerate}
  \item The defender observes the 24-dimensional state vector.
  \item An action is selected via $\varepsilon$-greedy exploration.
  \item The transition $(s, a, r, s', \text{done})$ is pushed to the replay buffer.
  \item If the buffer contains at least 64 transitions, one gradient update
        is performed with a randomly sampled mini-batch.
  \item Epsilon is decayed: $\varepsilon \leftarrow \max(0.05, 0.997\varepsilon)$.
  \item Every 150 steps, the target network is hard-copied from the policy network.
\end{enumerate}

Metrics recorded per episode: total reward, detection rate, false-positive
rate, average threat level, average training loss.
Training metrics are saved to \texttt{logs/metrics.csv} and visualised as
six-panel training curves. The DQN implementation and its training protocol
are unchanged from the prior version of this system; \S\ref{sec:dynamic_response}
explains why the DQN's own recorded results from this training run
(Chapter~\ref{ch:results}, Chapter~\ref{ch:discussion}) are used as direct
evidence when designing the dynamic-response mechanisms introduced there,
rather than retraining or extending the DQN itself.

%% ── 3.6 Gymnasium Environment ────────────────────────────────
\section{Gymnasium Environment (\texttt{CyberSecurityEnv})}
\label{sec:env}

\texttt{CyberSecurityEnv} extends \texttt{gymnasium.Env}~\citep{towers2024gymnasium}
and implements the standard Gymnasium interface.

\subsection{Observation Space}

The observation space is $\text{Box}(24,)$ with \texttt{float32} dtype.
The 24-dimensional vector is constructed as follows:
%
\begin{equation}
  s_t = \underbrace{[\,e_{a_t}^{(10)}\,]}_{\text{attack type (one-hot)}}
       \;\oplus\;
       \underbrace{[\,e_{k_t}^{(7)}\,]}_{\text{stage (one-hot)}}
       \;\oplus\;
       \underbrace{\left[\,T_t,\;\tfrac{n_t}{100},\;r_t\right]}_{\text{scalars}}
       \;\oplus\;
       \underbrace{[\,e_{\pi}^{(4)}\,]}_{\text{intent (one-hot)}},
  \label{eq:state}
\end{equation}
%
where $e_x^{(d)}$ denotes a $d$-dimensional one-hot vector for category $x$.
The scalar features ($T_t$, $n_t/100$, $r_t$) lie in $[0, 1]$ by
construction.
Including the attacker intent in the state vector allows both the DQN and the
SEDM to condition their responses on the inferred attacker type, even though
the SEDM queries the Markov chain directly.

This 24-dimensional layout, and the function that constructs it
(\texttt{encode\_state}), are shared verbatim between
\texttt{CyberSecurityEnv} and the live OpenCanary-integration pipeline
introduced in \S\ref{sec:dynamic_response}, so a state vector built from a
live event stream and one built from the synthetic training environment are
structurally identical, and everything validated against the synthetic
environment in this thesis applies unchanged to the live pipeline.

\subsection{Action Space}

The action space is \texttt{Discrete(5)}, one element per \texttt{HoneypotAction}.
The integer-to-action mapping is: 0\,=\,ALLOW, 1\,=\,LOG, 2\,=\,TROLL,
3\,=\,BLOCK, 4\,=\,ALERT.

\subsection{Episode Lifecycle}

\texttt{reset()} re-initialises the attacker Markov chains to their starting
states, draws a fresh per-session intensity/persona profile
(\S\ref{sec:realism}), clears the escalation history window and the EMA
accumulator (\S\ref{sec:escalation_tracking}), and resets the attack counter.
It returns the initial state vector and an \texttt{info} dictionary containing
the initial features and attack-type ground truth.

\texttt{step(action)} executes one simulation step:
%
\begin{enumerate}
  \item Calls \texttt{attacker.step()} to advance the Markov chains and
        sample a feature vector.
  \item Updates both escalation-tracking signals (\S\ref{sec:escalation_tracking}).
  \item Computes the threat level $T_t$ (Equation~\ref{eq:threat}), using
        whichever escalation signal is currently selected.
  \item Computes the reward $r_t$ (Section~\ref{sec:reward}).
  \item Constructs the next state vector $s_{t+1}$
        (Equation~\ref{eq:state}).
  \item Increments the step counter; terminates the episode if
        \texttt{max\_steps} is reached.
  \item Returns \texttt{(next\_state, reward, terminated, truncated, info)}.
\end{enumerate}

Episodes terminate by truncation (\texttt{max\_steps} reached) rather than
by natural termination conditions.
The environment does not model attacker defeat or withdrawal, consistent with
the assumption that the attacker's Markov chain is independent of the
defender's responses.

%% ── 3.7 Attack Classifier ────────────────────────────────────
\section{Random Forest Attack Classifier}
\label{sec:classifier}

The \texttt{AttackClassifier} wraps a \texttt{RandomForestClassifier} from
scikit-learn with a \texttt{StandardScaler} preprocessing step.

\subsection{Training Data Generation}

Training data is generated entirely from the parametric feature distributions
(Section~\ref{sec:features}), using the independent-sample override
described in \S\ref{sec:realism}.
For each of the ten attack types, 600 samples are generated independently:
%
\begin{equation}
  \mathcal{D}_{\text{train}} = \bigcup_{a=0}^{9}
  \bigl\{(\mathbf{x}, a) \,:\, \mathbf{x} \sim p(\cdot \mid a)\bigr\}_{600},
\end{equation}
%
yielding a balanced dataset of 6\,000 samples.
\texttt{class\_weight=`balanced'} provides a further guard against any
residual class imbalance.
The evaluation set (2\,000 samples, 200 per class) is generated with a
different random seed (999) to ensure independence from training data.

\subsection{Role in the System}

Two decision modes are evaluated (Chapter~\ref{ch:results}).
In \emph{oracle} mode, the SEDM reads the attack type directly from the
environment state vector (the ground-truth label) — this is an upper bound
on achievable performance, used to establish that the SEDM's decision logic
itself is correct given perfect information. In \emph{classifier-driven}
mode — the realistic operating condition, since ground truth is never
directly observable in deployment — the classifier's predicted attack type
is substituted into the state vector before the SEDM's matrix lookup, so
classification noise propagates into the policy's decisions. Both modes are
reported side by side in the extended evaluation (Chapter~\ref{ch:results});
classifier-driven decisions are used for the primary evaluation. The
classifier is retrained on the extended, session-coherent generator of
\S\ref{sec:realism} and re-evaluated under the identical protocol used to
produce the original held-out accuracy figure, so the two numbers reported
in Chapter~\ref{ch:results} are directly comparable.

%% ── 3.8 Evaluation Infrastructure ───────────────────────────
\section{Evaluation Infrastructure}
\label{sec:eval_infra}

\subsection{Metrics}

Two granularity levels are maintained:

\paragraph{Step-level (\texttt{StepRecord}).}
A dataclass capturing per-step data:
\texttt{episode, step, action, reward, attack\_type, kill\_chain\_stage,
threat\_level, is\_attack, predicted\_attack, loss, escalation\_rate}.

\paragraph{Episode-level (\texttt{EpisodeRecord}).}
Aggregated per-episode statistics:
%
\begin{align}
  \text{detection\_rate} &= \frac{\text{TP}}{\text{TP} + \text{FN}}, \\
  \text{false\_positive\_rate} &= \frac{\text{FP}}{\text{FP} + \text{TN}},
\end{align}
%
where a True Positive (TP) is any step where the attack was detected
(action $\neq$ ALLOW) and the traffic was genuinely malicious.
A False Positive (FP) is any step where the defender responded with a
non-ALLOW action on benign traffic.

\subsection{Experimental Protocol}

\paragraph{SEDM evaluation.}
The SEDM is evaluated for 30 episodes per attacker intent (120 total),
each episode lasting 200 steps, in the original protocol; Chapter~\ref{ch:results}
introduces a revised primary protocol using benign-traffic injection and a
larger episode budget, justified there directly against the statistical
evidence that motivates it.
No training is required; the policy is fully specified by the
Table~\ref{tab:sedm_matrix} and the Markov chain transition model.
Results are aggregated per intent and reported as mean $\pm$ standard
deviation over the evaluation episodes, with 95\% Wilson score confidence
intervals reported alongside point estimates for detection rate and
false-positive rate in the revised protocol.

\paragraph{DQN training.}
The DQN agent is trained for 300 episodes of 500 steps each using the
OPPORTUNISTIC intent.
All hyperparameters are listed in Table~\ref{tab:hyperparams}.
Training metrics are logged to \texttt{logs/metrics.csv}.

\paragraph{Reproducibility.}
All random seeds are fixed: seed 42 for training data generation and
environment resets; seed 999 for classifier evaluation.
The NumPy and PyTorch random states are seeded at the start of each run.
The new evaluation protocols introduced in \S\ref{sec:escalation_tracking}--\S\ref{sec:dynamic_response}
(escalation-mode comparison, reputation-crossing simulation, threshold-controller
convergence, and the revised primary protocol) follow the same seeding
discipline and are documented in full, with their exact procedures, in
Chapter~\ref{ch:results}.

%% ── 3.9 Escalation Tracking ──────────────────────────────────
\section{Escalation Tracking: Window and Severity-Weighted EMA}
\label{sec:escalation_tracking}

The escalation rate $r_t$ used throughout \S\ref{sec:threat} and
\S\ref{sec:sedm} was, in the prior version of this system, defined
exclusively as a hard sliding-window fraction:
%
\begin{equation}
  r_t^{\text{window}} = \frac{1}{|W|}\sum_{i \in W} \mathbb{1}[\text{attack}_i],
  \qquad |W| = 20.
\end{equation}
%
Two properties of this definition are worth making explicit, since they
motivate the alternative introduced here. First, it is \emph{binary per
step}: a Reconnaissance-severity probe and a Worms-severity outbreak
contribute identically to $r_t^{\text{window}}$, so the signal captures
attack \emph{frequency} but not attack \emph{severity}. Second, it has a
\emph{hard cutoff}: a step's contribution disappears the instant it exits
the trailing 20-step window, rather than fading gradually.

\paragraph{Severity-weighted exponential moving average.}
An alternative definition is introduced, computed alongside (not in place
of) the window signal:
%
\begin{equation}
  r_t^{\text{ema}} = \alpha \cdot s_{a_t} + (1-\alpha) \cdot r_{t-1}^{\text{ema}},
  \qquad \alpha = 0.15,
  \label{eq:ema}
\end{equation}
%
where $s_{a_t}$ is the same per-attack-type severity weight already defined
in Table~\ref{tab:attack_severity_meth}. This computation is $O(1)$ per
step, in contrast to the window signal's $O(|W|)$ deque maintenance; decays
smoothly rather than with a hard cutoff; and reflects the severity of
recent activity rather than merely its frequency. A direct consequence of
severity-weighting is that its steady state under sustained maximal-severity
activity is bounded above by the highest severity weight in the system
(0.90, for WORMS) rather than by 1.0 — a property that is not merely
theoretical but is confirmed empirically in Chapter~\ref{ch:results}, and
which has a direct implication for \texttt{RATE\_THRESHOLD} (R3,
\S\ref{sec:sedm_alg}): that threshold was calibrated against the window
signal's semantics, and its behaviour under the EMA signal is not
identical, motivating the independent empirical comparison in
Chapter~\ref{ch:results} before EMA tracking is relied upon operationally.

\paragraph{Deployment default.}
Both signals are always computed and exposed (as \texttt{escalation\_window\_rate}
and \texttt{escalation\_ema} respectively); an \texttt{escalation\_mode}
setting selects which one feeds $r_t$ in the equations of
\S\ref{sec:threat}--\S\ref{sec:sedm}. \texttt{"window"} remains the default
throughout this thesis except where explicitly stated otherwise, for the
calibration reason given above. The identical EMA computation is implemented
in the live OpenCanary-integration pipeline's session-tracking component
(\S\ref{sec:dynamic_response}), so the escalation signal is computed
consistently regardless of which of the system's two parallel data paths —
the synthetic training/evaluation environment, or a live event stream —
produced it.

%% ── 3.10 Auditable Dynamic Response ─────────────────────────
\section{Auditable Dynamic Response}
\label{sec:dynamic_response}

The mechanisms in \S\ref{sec:arch}--\S\ref{sec:escalation_tracking} are all
fixed at design time: once the SEDM table, its thresholds, and its override
rules are specified, none of them change in response to what the system
subsequently observes. This section introduces two mechanisms that do
respond to runtime observation, designed throughout to preserve the
auditability property that motivated the SEDM's original design — each
mechanism reduces, at any instant, to a small number of named numeric
quantities that a human analyst can read and verify directly, rather than
to a model whose behaviour must be probed indirectly.

\subsection{Cross-Session Reputation}
\label{sec:reputation}

\paragraph{Motivation.}
The evaluation environment and, in the live pipeline, the per-session state
tracker both discard session-scoped history once a session's time-to-live
has elapsed (nominally 300 seconds in the live pipeline). A source that
attacks, falls silent past that interval, and then returns is, under the
mechanisms described so far, indistinguishable from a source making genuine
first contact.

\paragraph{Mechanism.}
A \texttt{ReputationTracker} component maintains one offense score per
source identity, stored independently of the identity's individual
sessions, so it survives session expiry. The score decays continuously with
a configurable half-life $h$ rather than expiring at a fixed instant:
%
\begin{equation}
  \text{score}(t) = \text{score}(t_0) \cdot 0.5^{\,(t - t_0)/h}, \qquad
  h = 6 \text{ hours (default)}.
  \label{eq:reputation_decay}
\end{equation}
%
Each event increments the source's score by $\Delta \cdot s_a$ — an
\emph{offense increment} $\Delta = 0.25$ scaled by the event's attack-type
severity $s_a$ (Table~\ref{tab:attack_severity_meth}) — clamped to a
maximum of 1.0. This is deliberately the simplest possible update rule that
ties reputation growth to observed severity rather than to a flat per-event
increment: a NORMAL-labelled event contributes zero increment (since
$s_{\text{NORMAL}} = 0$), while sustained high-severity activity accumulates
reputation faster than sustained low-severity activity.

\paragraph{Integration with the SEDM.}
The resulting score feeds R4 (\S\ref{sec:sedm_alg}): once a source's decayed
reputation reaches \texttt{REPUTATION\_THRESHOLD} $= 0.60$, R4 fires ahead
of R1, upgrading the action by one level regardless of how benign the
immediate event appears. This ordering is a deliberate security-posture
decision, not an incidental consequence of implementation order: a source
that has already demonstrated sufficient reputation does not regain
unconditional trust merely by sending one innocuous-looking probe. The
explicit trade-off this implies — a shared or dynamically-reassigned
address that was briefly malicious remains penalised for the reputation
half-life even after genuinely malicious activity from that address has
ceased — is acknowledged here and revisited in Chapter~\ref{ch:discussion}.

\paragraph{Scope.}
\texttt{reputation} is threaded through the decision interface as an
explicit parameter (\texttt{MatrixPolicy.decide(..., reputation=...)}),
analogous to how \texttt{intent} is already threaded as an optional
argument, rather than being incorporated into the 24-dimensional state
vector itself: the state vector's dimensionality is load-bearing for the
classifier and the DQN's observation space (\S\ref{sec:env}), and no
persistent cross-episode source identity exists in the synthetic Gymnasium
environment in any case. The parameter defaults to 0.0 everywhere it is not
explicitly supplied, which reproduces the original three-rule override
behaviour exactly.

\subsection{Bounded Adaptation of the R3 Threshold}
\label{sec:threshold_controller}

\paragraph{Motivation and an explicit scope decision.}
Given the general aim of ``dynamic behavioural adjustment,'' the two
thresholds most immediately available for adaptation are
\texttt{ESC\_LOW\_THRESHOLD}/\texttt{ESC\_HIGH\_THRESHOLD} (Step 2,
\S\ref{sec:sedm_alg}) and \texttt{RATE\_THRESHOLD} (R3). This thesis adapts
only the second. $\rho(k, \pi)$, which the first two thresholds bucket, is
computed purely from the static, hand-specified Markov transition matrices
(\S\ref{sec:transition}) — it has no live observational component
whatsoever, so there is no runtime signal whose drift could honestly
indicate that the two band cuts require adjustment. \texttt{RATE\_THRESHOLD},
in contrast, gates on the escalation rate $r_t$, which \emph{is} a function
of observed traffic (\S\ref{sec:escalation_tracking}) — but the only
feedback genuinely available for it, in the absence of any analyst-labelling
pathway (Chapter~\ref{ch:discussion}), is R3's own trigger frequency, not
whether any individual trigger was the operationally correct decision. The
mechanism introduced here is accordingly scoped and named honestly: it
manages \textbf{alert-fatigue risk} (keeping the rate at which R3 fires
within an operator-chosen band, so that a downstream alert queue is neither
flooded nor allowed to fall silent), and makes no claim to improve
detection correctness, because no signal exists in this system that would
license such a claim.

\paragraph{Mechanism.}
\texttt{AdaptiveThresholds} implements a bounded deadband controller. Given
a target trigger rate $\hat{r}$, a tolerance $\tau$, a step size $\delta$,
and a bound $b$, the controller accumulates the R3 condition's truth value
over a rolling \texttt{observation\_window} of decisions (200 by default);
once the window is full, the current empirical trigger rate $\hat{r}_t$ is
compared against $[\hat{r} - \tau,\, \hat{r} + \tau]$ (defaults:
$\hat{r} = 0.10$, $\tau = 0.03$), and the threshold is nudged by $\pm\delta$
($\delta = 0.01$) in the corrective direction if the rate falls outside
that band, clamped throughout to $[\theta_0 - b,\, \theta_0 + b]$ around the
initial threshold $\theta_0$ (default bound $b = 0.10$). The R3 condition is
evaluated against the \emph{current} (possibly already-adjusted) threshold
before the R1/R2/R4 short-circuits in Step 4 of \S\ref{sec:sedm_alg}, and is
recorded by the controller on every decision regardless of whether R4, R1,
or R2 subsequently pre-empted R3 from actually firing.

\paragraph{Deployment default.}
\texttt{MatrixPolicy} accepts an optional \texttt{adaptive\_thresholds}
argument; its default value of \texttt{None} reproduces the static-threshold
behaviour of \S\ref{sec:sedm} exactly.

\subsection{Scope of ``Dynamic'' in This System}

Both mechanisms above satisfy a stricter standard than the phrase ``dynamic
behavioural adjustment'' might suggest by default: neither is a learned
function, and neither claims to improve decision \emph{correctness} beyond
what the static matrix and override rules of \S\ref{sec:sedm} already
provide. The reputation mechanism (\S\ref{sec:reputation}) changes
behaviour in response to a source's accumulated, decaying offense history —
a genuinely new capability relative to the system described in
\S\ref{sec:arch}--\S\ref{sec:eval_infra} — but its update rule is fixed
arithmetic. The threshold controller (\S\ref{sec:threshold_controller})
changes one numeric constant's value in response to its own trigger
frequency, within a bounded range, for an explicitly narrow, named purpose.
Neither mechanism requires training data, a reward signal, or a model whose
internal state cannot be read directly.

This scope was chosen deliberately, and Chapter~\ref{ch:discussion} makes
the argument explicit using this project's own DQN results
(\S\ref{sec:dqn_baseline}, Chapter~\ref{ch:results}) as direct supporting
evidence: a genuine analyst-feedback pathway — closing the loop from ``a
human confirmed this decision was right or wrong'' back into a trainable
signal — does not currently exist anywhere in this system, and building a
learning-based dynamic-adjustment mechanism without first building that
pathway would repeat, in miniature, the same failure mode this thesis's own
DQN baseline already exhibited: optimising against a distributional proxy
in the absence of genuine ground truth. Chapter~\ref{ch:discussion}
identifies constructing that pathway as the specific prerequisite that
would need to exist before this scope decision should be revisited.
